\section{Metaphysical Mutations}

\begin{quotationx}
Metaphysical mutations -- that is to say radical, global transformations in the values to which the majority subscribe -- are rare in the history of humanity. The rise of Christianity might be cited as an example.

Once a metaphysical mutation has arisen, it tends to move inexorably toward its logical conclusion. Heedlessly, it sweeps away economic and political systems, aesthetic judgments and social hierarchies. No human agency can halt its progress -- nothing except another metaphysical mutation.

It is a fallacy that such metaphysical mutations gain ground only in weakened or declining societies. When Christianity appeared, the Roman Empire was at the height of its powers: supremely organized, it dominated the known world; its technological and military prowess had no rival. Nevertheless, it had no chance. When modern science appeared, medieval Christianity was a complete, comprehensive system which explained both man and the universe; it was the basis for government, the inspiration for knowledge and art, the arbiter of war as of peace and the power behind the production and distribution or wealth --- none of which was sufficient to prevent its downfall.


\flright{\textsc{Michel Houellebecq}, \emph{Elementary Particles}}
\end{quotationx}



\flrightit{Posted on 2008-09-14 by Cologero}

